% ============================================
% Johns Hopkins Letter Template
% ============================================

\documentclass[11pt,letterpaper]{letter}

\usepackage{graphicx}
\usepackage{eso-pic}
\usepackage{geometry}
\usepackage{microtype}
\usepackage[utf8]{inputenc}
\usepackage[hidelinks]{hyperref}

\geometry{left=25mm,right=25mm,top=25mm,bottom=25mm}

\newcommand{\PutJHULogoFG}[1]{%
  \AddToShipoutPictureFG*{%
    \AtPageUpperLeft{%
      \hspace{10mm}%
      \raisebox{-38mm}{%
        \includegraphics[width=6cm]{#1}%
      }%
    }%
  }%
}

\signature{Decory Edwards}
\address{
Department of Economics \\
Johns Hopkins University \\
Baltimore, MD 21218 \\
\texttt{dedwar65@jh.edu}
}

\begin{document}
\PutJHULogoFG{Images/jhulogo.png}

\begin{letter}{
Hiring Committee \\
Department of Economics \\
University of California, Los Angeles
}

\opening{Dear Hiring Committee,}

I am writing to express my interest in the Adjunct Professor position in the Department of Economics at UCLA. I am a Ph.D.\ candidate in the Department of Economics at Johns Hopkins University (advisors: Christopher D.\ Carroll, Nicholas W.\ Papageorge, and Stelios Fourakis), and I expect to receive my degree in May 2026. My primary specializations are macroeconomic modeling, wealth and income dynamics, and computational economics.

My research agenda focuses on the ability of macroeconomic models to generate empirically realistic distributions of wealth and income over the life cycle. A key driver of that work is modeling heterogeneity in the rate of return on assets, and understanding how return dispersion contributes to portfolio accumulation across households.

In my job market paper, I calibrate a life cycle heterogeneous agent model to match wealth moments from the Survey of Consumer Finances by allowing households to differ in their portfolio returns. The model helps explain observed wealth concentration and return dispersion. In a related research project using the Health and Retirement Study, I examine how trust influences both income growth and asset returns. I find a hump shaped relationship for trust and returns and characterize the distribution of the persistent component of returns to net worth. These experiences have equipped me to present technical material clearly, integrate quantitative tools into instruction, and communicate rigorous economic reasoning to students at varying levels of preparation.

I am particularly drawn to UCLA’s Department of Economics because of its strong undergraduate program and its emphasis on high quality, research informed teaching in large class environments. I am prepared to teach Principles of Economics, Intermediate Microeconomics, and Intermediate Macroeconomics, as well as Statistics and Econometrics. My training in dynamic optimization, econometric modeling, and computational methods allows me to teach econometrics with both mathematical rigor and practical application. I regularly use Python and Stata for data analysis and simulation, and I incorporate coding demonstrations, empirical replication exercises, and structured problem solving into my teaching. I am also enthusiastic about employing online tools, interactive assignments, and digital learning platforms to enhance engagement in larger lecture settings. My goal is to make quantitative material accessible without sacrificing analytical depth.

I am enthusiastic about living and working in Los Angeles and participating fully in UCLA’s in person academic community. I value the energy and intellectual diversity of a large public research university and would welcome the opportunity to contribute through teaching, service, and continued scholarly engagement. I am also attracted to the possibility of a multi year adjunct appointment, which would allow me to refine large scale undergraduate instruction while continuing to develop my broader research agenda within a vibrant academic environment.

I believe I would be a strong fit for this role because:
\begin{enumerate}
    \item I have strong preparation to teach statistics, econometrics, and core theory courses with quantitative rigor.
    \item I am comfortable teaching large classes and using technology to support student learning.
    \item I bring a research active perspective that enriches undergraduate instruction while aligning with UCLA’s standards of excellence.
\end{enumerate}

Thank you very much for your time and consideration. I would welcome the opportunity to contribute to UCLA’s Department of Economics.

\closing{Sincerely,}

\end{letter}
\end{document}
